%; whizzy paragraph -pdf xpdf -latex ./whizzypdfptex.sh
%; whizzy-paragraph "^\\\\begin{frame}"
% latex beamer presentation.
% platex, latex-beamer でコンパイルすることを想定。 

%     Tokyo Debian Meeting resources
%     Copyright (C) 2009 Junichi Uekawa
%     Copyright (C) 2009 Nobuhiro Iwamatsu

%     This program is free software; you can redistribute it and/or modify
%     it under the terms of the GNU General Public License as published by
%     the Free Software Foundation; either version 2 of the License, or
%     (at your option) any later version.

%     This program is distributed in the hope that it will be useful,
%     but WITHOUT ANY WARRANTY; without even the implied warreanty of
%     MERCHANTABILITY or FITNESS FOR A PARTICULAR PURPOSE.  See the
%     GNU General Public License for more details.

%     You should have received a copy of the GNU General Public License
%     along with this program; if not, write to the Free Software
%     Foundation, Inc., 51 Franklin St, Fifth Floor, Boston, MA  02110-1301 USA

\documentclass[cjk,dvipdfmx,12pt]{beamer}
\usetheme{Tokyo}
\usepackage{monthlypresentation}

%  preview (shell-command (concat "evince " (replace-regexp-in-string "tex$" "pdf"(buffer-file-name)) "&")) 
%  presentation (shell-command (concat "xpdf -fullscreen " (replace-regexp-in-string "tex$" "pdf"(buffer-file-name)) "&"))
%  presentation (shell-command (concat "evince " (replace-regexp-in-string "tex$" "pdf"(buffer-file-name)) "&"))

%http://www.naney.org/diki/dk/hyperref.html
%日本語EUC系環境の時
\AtBeginDvi{\special{pdf:tounicode EUC-UCS2}}
%シフトJIS系環境の時
%\AtBeginDvi{\special{pdf:tounicode 90ms-RKSJ-UCS2}}

\title{東京エリアDebian勉強会}
\subtitle{第75回 2011年4月度}
\author{まえだこうへい mkouhei@debian.org\\IRC nick: mkouhei}
\date{2011年4月16日}
\logo{\includegraphics[width=8cm]{image200607/openlogo-light.eps}}

\begin{document}

\frame{\titlepage{}}

\section{}
\begin{frame}
 \frametitle{Agenda}
\begin{minipage}[t]{0.45\hsize}
  \begin{itemize}
  \item 注意事項
	\begin{itemize}
	 \item 飲酒禁止
	 \item 宗教禁止
	 \item 営利活動禁止
	\end{itemize}
   \item 最近あったDebian関連のイベント報告
	\begin{itemize}
	 \item OSC 2011 Tokyo/Spring
	\end{itemize}
 \end{itemize}
\end{minipage} 
\begin{minipage}[t]{0.45\hsize}
 \begin{itemize}
  \item 会長就任挨拶
  \item backports.debian.orgの話
  \item initramfs-toolsの話
  \item 月刊PPC64ポーティング
  \item 僕がDD目指すの手伝ってください
 \end{itemize}
\end{minipage}
\end{frame}

\begin{frame}
 \frametitle{前回}
\begin{minipage}[t]{0.45\hsize}
  \begin{itemize}
  \item 注意事項
	\begin{itemize}
	 \item 飲食禁止
	 \item 宗教禁止
	 \item 営利活動禁止
	\end{itemize}
  \end{itemize}
\end{minipage}
\begin{minipage}[t]{0.45\hsize}
\begin{itemize}
  \item 最近あったDebian関連のイベント報告
	\begin{itemize}
	 \item CACert認証会 \& GPGキーサイン
	 \item CACert公式トレーニング
	\end{itemize}
 \end{itemize}
\end{minipage}
\end{frame}


\emtext{イベント報告}

\section{DWN quiz}
\begin{frame}{Debian 常識クイズ}

Debian の常識、もちろん知ってますよね?
知らないなんて恥ずかしくて、知らないとは言えないあんなことやこんなこと、
みんなで確認してみましょう。

今回の出題範囲は\url{debian-devel-announce@lists.debian.org} に投稿された
内容とDebian Project Newsからです。

\end{frame}

\subsection{問題}
%; whizzy-master ../debianmeetingresume201101.tex
% $B0J>e$N@_Dj$r$7$F$$$k$?$a!"$3$N%U%!%$%k$G(B M-x whizzytex $B$9$k$H!"(Bwhizzytex$B$,MxMQ$G$-$^$9!#(B
%
% $B$A$J$_$K!"%/%$%:$OJL%V%i%s%A$G:n@.$7!"$N$A$K%^!<%8$7$^$9!#5U$K%^!<%8$7(B
% $B$J$$$h$&$K$7$^$7$g$&!#(B
% (shell-command "git checkout quiz-prepare")

\santaku
{2011$BG/EY(B Debian JP $B2qD9$OC/$G$7$g$&$+!)(B}
{$BA0ED(B $B9LJ?(B}
{$B4d>>(B $B?.MN(B}
{$B9SLZ(B $BLw9((B}
{A}
{}

\santaku
{2011$BG/EY(B DPL $B$OC/$G$7$g$&$+!)(B}
{Kurt Roeckx}
{Kenshi Muto}
{Stefano Zacchiroli}
{C}
{}

\santaku
{Debian Policy 3.9.2 $B$GDI2C$5$l$F$$$J$$9`L\$O$I$l$+(B}
{Debian$B%"%+%&%s%H$r%a%s%F%J%"%I%l%9$KDI2C$9$kI,MW$,$"$j$^$9!#(B}
{$B%"!<%-%F%/%A%c0MB8$N%i%$%V%i%jEy$O(B DEB\_HOST\_MULTIARCH $B$G<hF@$7$?CM$rMx(B
$BMQ$9$kI,MW$,$"$j$^$9!#(B}
{$BA4$F$N%Q%C%1!<%8$O(BVCS$B$G4IM}$9$kI,MW$,$"$j$^$9!#(B}
{C}
{}

\santaku
{DebConf chairs$B$K;XL>$5$l$?$N$OC/!)(B}
{Gunnar Wolf}
{Maeda Kouhei}
{Junichi Uekawa}
{A}
{}

\santaku
{ries.debian.org $B$G<hF@$G$-$k$h$&$K$J$C$?%G!<%?$O!)(B}
{debian.org $B$N2TF/>u67%G!<%?(B}
{debian.org $B$N(BML$B%G!<%?$r(Bgzip$B$G8G$a$?J*(B}
{dak $B$N%G!<%?(B}
{C}
{}

\santaku
{/run $B$,>C$5$l$?M}M3$O!)(B}
{/go $B$N$[$&$,$h$/$M!)$H$$$&?M$,8=$l$?!#(B}
{initscripts $B$,$^$@%5%]!<%H$7$F$J$$!*(B}
{/run $B$rF~$l$?$N$O%Q%C%1!<%8%s%0%_%9$G$9!#(B}
{B}
{}


\emtext{会長就任挨拶}

\emtext{prework}

{\footnotesize
 %; whizzy-master ../debianmeetingresume201101.tex
% $B0J>e$N@_Dj$r$7$F$$$k$?$a!"$3$N%U%!%$%k$G(B M-x whizzytex $B$9$k$H!"(Bwhizzytex$B$,MxMQ$G$-$^$9!#(B


\begin{prework}{ $B5HLn(B(yy\_y\_ja\_jp) }

$BIaCJDL$j%G%9%/%H%C%W4D6-$H$7$F;H$C$F$^$9!%(B
\end{prework}

\begin{prework}{ henrich }

$B3hMQ$OFCCJ$7$F$$$^$;$s!#DdEE$N$?$a$K2H%5!<%P$HCN?M$N2q<R$KCV$$$F$"$k%5!<%P$NN>J}$,;_$^$C$FFq57$7$^$7$?!#(B
\end{prework}

\begin{prework}{ dictoss($B?yK\!!E5=<(B) }

 \begin{itemize}
  \item $BDdEE$G(BPC$B$r5/F0$G$-$J$$$H$-$b$"$j!"$`$7$m(BDebian$B$r;H$($J$$$H$-$b$"$j$^$7$?!#(B
  \item $B;H$$J}$O%a!<%k$7$?$j!"(BWeb$B%V%i%&%:$r$7$?$j$HIaCJ$HJQ$o$i$J$$46$8$G(B
	$B$9!#(B
 \end{itemize}
\end{prework}

\begin{prework}{ $B4d>>(B $B?.MN(B }

 \begin{itemize}
  \item chromium $B$N3HD%%W%m%0%i%`$r=q$$$F!"CO?L$,Mh$?$i%&%#%s%I%&$r=P$9$h$&$J%W%m%0%i%`$r=q$$$F$_$^$7$?!#(B
  \item $BDdEEMQ$K(BVPS$B$r7@Ls$7$F!"(BDebian$B$K$7$?!#(B
  \item Debian$B>e$G(BOSM$B$r$$$8$C$?$j$7$?!#(B
 \end{itemize}
 
\end{prework}

\begin{prework}{ $BLnEg!!5.1Q(B }

$B?L:R$GBgJQ$JCf!"Hs>o$K$D$^$i$J$$Nc$GCQ$:$+$7$$$N$G$9$,!"(B

 \begin{enumerate}
  \item $B?L:R>pJs=8$a$O<g$K(Bdebain+ephipany$B$NAH$_9g$o$;$,B?$$$G$9(B
  \item $B?L:R$K4X$9$k=t!9$N>pJs=8$a$G(BNIKKEINET($BM-NA!K$rMxMQ$9$k:]$O(Bdebian+iceweasel$B$P$+$j$G$9!#(B
  \item USB$B%V!<%H2DG=$J%Q%=%3%s$5$($"$l$PM-;v$N:]$I$3$G$bBP1~$G$-$k$h$&(B
	$B$K<+J,$N(Bdebian$B4D6-$rF~$l$?%V!<%H2DG=$J(BUSB$B%a%b%j$r;}$AJb$$$F$$$?(B
	$B$j$7$^$9!#<+M3$J;H$$J}$G$-$k<+J,$N4D6-$r%]%1%C%H$K$7$N$P$;$F;}$A(B
	$BJb$1$k$N$O$h$$$+$b!#CO?L$"$C$?$H$-$O$3$N(BUSB$B$r6;%]%1%C%H$KF~$l$F(B
	$BHrFq$7$^$7$?!#!J$?$^$?$^F~$C$F$$$?$H$b$$$&!#$G$b<+J,$N4D6-$O$3$3(B
	$B$K$"$k$+$i$H$A$g$C$H0B?4$G$-$?$N$O3N$+!#!K(B
 \end{enumerate}
 $B"(JL$K(Bdebian$B$G$J$/$F$b$h$$$8$c$s$H$$$&%D%C%3%_$O%J%7$NJ}8~$G!#(B
\end{prework}

\begin{prework}{ $B;3ED(B }

$B?L:R4X78$@$H$"$^$j3hLv$7$F$$$J$$$G$9!D(B

$B:#2s$O(BPC$B$r;}$C$FHo:R$7$?$b$N$N$[$H$s$I=PHV$,$J$/!"(B
$B7HBSEEOC$@$1$,3hLv$7$^$7$?!#(B

$B$=$&$$$&$o$1$G!"Aa5^$K(BAndroid$B7HBS$K(BDebian/Ubuntu$B$r(B
$BF~$l$?$$$H;W$$$^$9!J(BiDroid$B$GDLOC$G$-$k$N$@$m$&$+!)!K!#(B
\end{prework}

\begin{prework}{ $BB<ED?.?M(B }

$B?L:R$KBP$7$FD>@\$O3hMQ$G$-$F$$$^$;$s!#$?$@!"(BSahana Eden$B$NF|K\%A!<%`$N(BLaunchpad, Bazaar$B$KBP$9$k5?Ld$KEz$($k$3$H$G4V@\E*$K9W8%$G$-$l$P$H!#(B
\end{prework}

\begin{prework}{ Kiwamu Okabe }

sakura VPS$B>e$K(BDebian$B$r%$%s%9%H!<%k$7$F$$$^$9!#(B
$B;E;v$G(BNetBSD$B$r;H$C$?3+H/$r$7$F$$$k$N$G$9$,%M%C%H%o!<%/$,;_$^$C$F$b%a%s%P!<$,3+H/7QB3$G$-$k$h$&$K8x<0(Bcvs$B%j%]%8%H%j$N(Btar$B6L$r<h$k$h$&$K$7$^$7$?!#(B

netbsd\_cvs =rsync=$>$ sakura\_vps(debian) =tar$B6L(B=$>$ $B2q<R$N%a%s%D(BPC

http://www.masterq.net/files/
$BKhD+(B4$B;~$4$m99?7$7$F$$$^$9!#(B
\end{prework}

\begin{prework}{ Osamu MATSUMOTO }

$BL$Dj$@$C$?$N$G2sEz$J$7$G!#(B
$B$b$&(B10$BG/$0$i$$(BDebian$B;H$C$F$P$+$j$J$N$G!"9W8%$7$?$$$H;W$C$F$$$k!#(B
\end{prework}

\begin{prework}{ yamamoto }

$B:#2s$N?L:R$G5"Bp;~4V$,Aa$/$J$C$?J,!"%^%7%s$X$N$*?($j$,A}$($^$7$?!#(B
$B:#$b$*$$$i$N2#$G%&%s%&%sS9$C$F$$$^$9!#(B
\end{prework}

\begin{prework}{ $B$^$($@$3$&$X$$(B }

$B2f$,2H$N(BDebian$B%^%7%s$?$A!#(B
\begin{itemize}
 \item $BCO?L$NEvF|!"%h%a$H%G%#%:%K!<%7!<$+$i5"$l$J$/!"MbF|5"Bp$9$k$^$G(B
       $B$3$^$a$,?4G[$G;EJ}$J$+$C$?$N$G!"(BOpenBlockS 600$B$K(BWeb$B%+%a(B
       $B%i$r$D$1$F!"$3$^$A$c$s4F;k%7%9%F%`$r:n$j$^$7$?!#IaCJ$NMM;R$d!"Dd(B
       $BEE$N3NG'!"Bg$-$JM>?L$,$"$C$?:]$N<+Bp$NMM;R$r8+$F$^$9!#(B
 \item $B0lJ}!"(BOpenBlockS 266$B$G%O!<%I%G%#%9%/$G2TF/$7$F$$$?IaDL$N%5!<%P$?$A(B(Web,
       DNS, CouchDB, Git$B%j%b!<%H%j%]%8%H%jMQ(B, $B4F;k$J$I(B)$B$O!"DdEE$KHw$($9(B
       $B$Y$FDd;_$7$^$7$?!#(B
 \item $BCO?L$ND>8e$O$7$P$i$/(BMacBook$B$N;HMQ$O95$($F!"(BiPad$B$d(BAndroid$B;H$C$F$^(B
       $B$7$?!#(B
 \item $B:#8e$O30$K=P$;$k%5!<%P$d%G!<%?$O!"4pK\%/%i%&%I%5!<%S%9$rMxMQ$7$F!"(B
       OpenBlockS$BEy$O%G%#%9%/%l%92=$7$F!"%;%s%5!<$G$N4F;k@lMQ$K$7$F$$$/(B
       $BM=Dj$G$9!#(B
\end{itemize}

\end{prework}

}


\emtext{backports.debian.orgの話}

\emtext{initramfs-toolsの話}

\emtext{月刊PPC64ポーティング}

\emtext{僕がDD目指すの手伝ってください}

\begin{frame}{今後のイベント}
 
\begin{itemize}
 \item 5月 OSC仙台
 \item 6月　Debian勉強会
 \item 7月　Debian勉強会 \& Debconf11 in サラエボ
\end{itemize}
未実施ネタ：PS Moveネタ? デジタル放送取り込み? Debian Pod cast? 100台Squeezeアップグレード(吐血)体験記? 
\end{frame}

\begin{frame}{今日の宴会場所}

神保町周辺。

\end{frame}

\end{document}

;;; Local Variables: ***
;;; outline-regexp: "\\([ 	]*\\\\\\(documentstyle\\|documentclass\\|emtext\\|section\\|begin{frame}\\)\\*?[ 	]*[[{]\\|[]+\\)" ***
;;; End: ***
